\documentclass[11pt]{article}

% ============================================================================
%  The Peierls Argument: a complete proof
%
%  Companion document to Ising_Model.tex (the structured transcription of
%  Chapter 3 of Friedli & Velenik, "Statistical Mechanics of Lattice
%  Systems"). NOTE: this entire document is an ADDITION beyond the book's
%  transcription: it formalizes in full the proofs that Ising_Model.tex only
%  summarizes in Section 3.7.2 (Lemma 3.37, an entropy bound replacing
%  Lemma 3.38, and the deduction of a phase transition at low temperature,
%  item 3 of Theorem 3.25).
%  References of the form "Lemma 3.xx" point to the numbering of the
%  transcription (= the book); internal results are numbered independently.
% ============================================================================

\usepackage[utf8]{inputenc}
\usepackage[T1]{fontenc}
\usepackage{amsmath,amssymb,amsthm,mathtools,mathrsfs}
\usepackage[margin=1in]{geometry}
\usepackage{hyperref}

% --- theorem environments (as in Ising_Model.tex) ---------------------------
\theoremstyle{plain}
\newtheorem{theorem}{Theorem}[section]
\newtheorem{lemma}[theorem]{Lemma}
\newtheorem{proposition}[theorem]{Proposition}
\newtheorem{corollary}[theorem]{Corollary}
\theoremstyle{definition}
\newtheorem{definition}[theorem]{Definition}
\theoremstyle{remark}
\newtheorem{remark}[theorem]{Remark}

% --- notation: identical to Ising_Model.tex ---------------------------------
\newcommand{\Z}{\mathbb{Z}}
\newcommand{\R}{\mathbb{R}}
\newcommand{\Rgeq}{\R_{\ge 0}}
\newcommand{\fin}{\Subset}          % finite subset of Z^d
\newcommand{\incr}{\Uparrow}        % van Hove convergence to Z^d
\newcommand{\La}{\Lambda}
\newcommand{\sg}{\sigma}
\newcommand{\Ham}{\mathcal{H}}
\newcommand{\Eset}{\mathscr{E}}
\newcommand{\avg}[1]{\langle #1 \rangle}
\newcommand{\meanp}[1]{\langle #1 \rangle^{+}}
\newcommand{\meanm}[1]{\langle #1 \rangle^{-}}
% --- additional notation for this document only -----------------------------
\newcommand{\dualZ}{(\Z^2)^*}       % dual lattice Z^2 + (1/2,1/2)
\newcommand{\Intr}{\mathrm{Int}}    % interior of a contour
\newcommand{\Extr}{\mathrm{Ext}}    % exterior of a contour
\newcommand{\Gam}{\Gamma}           % contour family of a configuration
\newcommand{\dedges}{\mathscr{E}^*} % dual disagreement edges

\title{\textbf{The Peierls Argument: a complete proof}\\[4pt]
\large Companion to the transcription of Chapter 3 of\\
Friedli \& Velenik, \emph{Statistical Mechanics of Lattice Systems}}
\author{}
\date{}

\begin{document}
\maketitle

\noindent\fbox{\parbox{\dimexpr\linewidth-2\fboxsep-2\fboxrule}{%
\textbf{Status of this document.} Everything below is an \emph{addition}
relative to the transcription \texttt{Ising\_Model.tex}: it expands the proof
summaries of \S3.7.2 there into complete proofs. Statements quoted from the
book keep their book numbering in brackets; all other formulations,
intermediate lemmas and proof details are ours and should not be mistaken for
the source text.}}

\section{Setting and main statement}
\label{sec:setting}

Throughout, $d=2$ and $h=0$. We work in the boxes
$B(n)=\{-n,\dots,n\}^2\fin\Z^2$ with the $+$ boundary condition: by
Definition~3.3 of the transcription, the Gibbs distribution
$\mu^{+}_{B(n);\beta,0}$ on
\[
  \Omega^{+}_{B(n)}
  =\{\omega\in\{-1,1\}^{\Z^2} : \omega_i=+1 \ \forall i\notin B(n)\}
\]
is
\[
  \mu^{+}_{B(n);\beta,0}(\omega)
    =\frac{e^{-\Ham^{+}_{B(n);\beta,0}(\omega)}}{Z^{+}_{B(n);\beta,0}},
  \qquad
  \Ham^{+}_{B(n);\beta,0}(\omega)
    =-\beta\!\!\sum_{\substack{\{i,j\}\in\Eset_{\Z^2}\\ \{i,j\}\cap B(n)\neq\varnothing}}\!\!
       \sg_i(\omega)\sg_j(\omega),
\]
where $\Eset_{\Z^2}$ denotes the set of nearest-neighbour edges of $\Z^2$ and
spins outside $B(n)$ are frozen to $+1$.

The goal of this document is a complete proof of:

\begin{theorem}[low-temperature phase transition; existence part of item 3 of Theorem 3.25]
\label{thm:main}
Let $d=2$, $h=0$. There exists $\beta_0<\infty$ such that for every
$\beta\ge\beta_0$,
\[
  m^*(\beta)=\meanp{\sg_0}_{\beta,0}\ \ge\ \tfrac12\ >\ 0 .
\]
Consequently $\meanp{\cdot}_{\beta,0}\neq\meanm{\cdot}_{\beta,0}$ for all
$\beta\ge\beta_0$ (non-uniqueness, Definition~3.24), and
$\beta_c(2)\le\beta_0<\infty$. One may take $\beta_0=1$.
\end{theorem}

The proof occupies the rest of the document: contours are constructed in
Section~\ref{sec:contours}, Peierls' estimate (Lemma~3.37 of the
transcription) is proved in Section~\ref{sec:peierls}, an entropy bound
(replacing Lemma~3.38 there by a simpler walk count, sufficient for our
purposes) in Section~\ref{sec:counting}, and the conclusion is drawn
in Section~\ref{sec:conclusion}.

\section{Contours}
\label{sec:contours}

Throughout this section $n$ is fixed and $\omega$ denotes a configuration in
$\Omega^{+}_{B(n)}$. We construct the contour representation of $\omega$
underlying \S3.7.2 of the transcription, and we establish the geometric and
combinatorial lemmas used in
Sections~\ref{sec:peierls}--\ref{sec:conclusion}. Everything is elementary;
the only two facts we use without proof are the Jordan curve theorem for
polygonal curves and its standard crossing-parity refinement, both flagged
explicitly below.

\subsection{The dual lattice and disagreement edges}

The \emph{dual lattice} is
\[
  \dualZ \;=\; \Z^2+\big(\tfrac12,\tfrac12\big)
  \;=\;\big\{\,(a+\tfrac12,\,b+\tfrac12) : a,b\in\Z \,\big\},
\]
whose vertices are the centres of the unit faces (plaquettes) of $\Z^2$; its
edges are the unit segments joining dual vertices at Euclidean distance $1$,
so that each dual vertex has exactly four incident dual edges. Every
nearest-neighbour edge $e=\{i,j\}\in\Eset_{\Z^2}$ determines a unique
\emph{dual edge} $e^*$, namely the edge of $\dualZ$ crossing $e$
perpendicularly at its midpoint:
\begin{align*}
  \{(a,b),(a+1,b)\}^*&=\big\{(a+\tfrac12,b-\tfrac12),\,(a+\tfrac12,b+\tfrac12)\big\},\\
  \{(a,b),(a,b+1)\}^*&=\big\{(a-\tfrac12,b+\tfrac12),\,(a+\tfrac12,b+\tfrac12)\big\},
\end{align*}
and $e\mapsto e^*$ is a bijection from $\Eset_{\Z^2}$ onto the set of edges
of $\dualZ$. Given $\omega\in\Omega^{+}_{B(n)}$, define the set of
\emph{disagreement edges}
\[
  \dedges(\omega)
  \;=\;\big\{\, e^* \;:\; e=\{i,j\}\in\Eset_{\Z^2},\
        \sg_i(\omega)\sg_j(\omega)=-1 \,\big\}.
\]
This set is finite: a disagreeing pair $\{i,j\}$ has at least one endpoint
with spin $-1$, hence in $B(n)$ (spins outside $B(n)$ are $+1$), and each
site of $B(n)$ has four incident edges, so
$|\dedges(\omega)|\le 4|B(n)|<\infty$.

The construction of contours rests on the following parity property.

\begin{lemma}[even-degree property]
For every $\omega\in\Omega^{+}_{B(n)}$ and every dual vertex $x^*\in\dualZ$,
the number of edges of $\dedges(\omega)$ incident to $x^*$ is $0$, $2$
or~$4$.
\end{lemma}

\begin{proof}
Write $x^*=(a+\tfrac12,b+\tfrac12)$ and let $i_1=(a,b)$, $i_2=(a+1,b)$,
$i_3=(a+1,b+1)$, $i_4=(a,b+1)$ be the four corners of the plaquette centred
at $x^*$. From the explicit description of dual edges above, the four edges
of $\dualZ$ incident to $x^*$ are precisely the duals of the four plaquette
edges $\{i_1,i_2\}$, $\{i_2,i_3\}$, $\{i_3,i_4\}$, $\{i_4,i_1\}$. With
indices taken mod $4$,
\[
  \prod_{m=1}^{4}\sg_{i_m}(\omega)\,\sg_{i_{m+1}}(\omega)
  \;=\;\prod_{m=1}^{4}\sg_{i_m}(\omega)^2\;=\;+1,
\]
since every spin occurs in exactly two factors. Each factor equals $\pm1$,
so the number of indices $m$ with $\sg_{i_m}\sg_{i_{m+1}}=-1$ --- which is
the degree of $x^*$ in $\dedges(\omega)$ --- is even.
\end{proof}

\subsection{The deformation rule and the contour decomposition}

For a dual vertex $x^*$, call the four incident dual edges its \emph{north},
\emph{east}, \emph{south} and \emph{west} edges (joining $x^*$ to
$x^*+(0,1)$, $x^*+(1,0)$, $x^*-(0,1)$, $x^*-(1,0)$ respectively).

\medskip\noindent\textbf{Deformation rule.}
At every dual vertex $x^*$ incident to at least one edge of
$\dedges(\omega)$, partition the incident edges of $\dedges(\omega)$ into
pairs as follows (this is possible by the even-degree property just proved):
\begin{itemize}
  \item if exactly two edges of $\dedges(\omega)$ are incident to $x^*$,
        pair them with each other;
  \item if all four edges incident to $x^*$ belong to $\dedges(\omega)$,
        pair the north edge with the east edge, and the south edge with the
        west edge.
\end{itemize}
Each pair at $x^*$ is thought of as a \emph{passage} of a curve through
$x^*$: a \emph{corner} if the two edges are perpendicular, \emph{straight}
if they are collinear. At a vertex of degree $4$ the rule produces two
corners (the north--east and the south--west corner) which do not cross each
other; the symmetric convention (north--west with south--east) would serve
equally well, whereas pairing north with south and east with west is
excluded, as it would produce two straight passages crossing at $x^*$.
Nothing below depends on which of the two admissible conventions is fixed;
we fix the north--east/south--west one once and for all.

\begin{lemma}\label{lem:contour-decomp}
Let $\omega\in\Omega^{+}_{B(n)}$ with $\dedges(\omega)\neq\varnothing$.
There is a unique partition
\[
  \dedges(\omega)\;=\;\gamma_1^*\sqcup\cdots\sqcup\gamma_k^*
\]
such that each $\gamma_m^*$ is the edge set of a closed circuit of $\dualZ$
which traverses every edge of $\gamma_m^*$ exactly once
(edge-self-avoiding) and which moves through each visited dual vertex along
a pair prescribed by the deformation rule (i.e.\ any two consecutive edges
of the circuit are paired at their common endpoint). Moreover, every
$|\gamma_m^*|$ is even and $\ge 4$.
\end{lemma}

\begin{proof}
For $e^*\in\dedges(\omega)$ and an endpoint $x^*$ of $e^*$, the deformation
rule assigns to $e^*$ a unique \emph{partner} at $x^*$: the edge paired with
$e^*$ among the edges of $\dedges(\omega)$ incident to $x^*$. The partner
relation at $x^*$ is a fixed-point-free involution of those edges. Let $L$
be the graph with vertex set $\dedges(\omega)$ in which $e^*$ and $f^*$ are
adjacent when they are partners at a common endpoint.

Each $e^*$ has exactly one partner at each of its two endpoints, so exactly
two partnership links, attached at \emph{distinct} endpoints of $e^*$.
The two partners are moreover distinct edges: if $f^*$ were the partner of
$e^*$ at both endpoints of $e^*$, then $f^*\neq e^*$ would be a dual edge
with the same two endpoints as $e^*$, impossible since $\dualZ$ has no
multiple edges. Hence $L$ is $2$-regular, with no loops and no double
adjacencies, and is therefore a disjoint union of cycles
$(e_1^*,e_2^*,\dots,e_r^*)$ with $r\ge3$, consecutive elements (cyclically)
being partners at a common endpoint.

Fix such a cycle and let $x_m^*$ denote the common endpoint at which
$e_m^*$ and $e_{m+1}^*$ are partners (indices mod $r$). Since the two links
of $e_m^*$ are attached at its two distinct endpoints, the endpoints of
$e_m^*$ are exactly $x_{m-1}^*$ and $x_m^*$. Thus
$x_0^*,x_1^*,\dots,x_{r-1}^*,x_0^*$ is a closed walk in $\dualZ$ which
traverses $e_1^*,\dots,e_r^*$, each exactly once, and which moves through
every visited vertex along a prescribed pair. The edge sets of the cycles
of $L$ therefore partition $\dedges(\omega)$ into circuits as in the
statement.

\emph{Uniqueness.} In any partition as in the statement, when a circuit
traverses an edge $e^*$, the preceding edge is the partner of $e^*$ at the
entry endpoint and the following edge is the partner at the exit endpoint;
hence the edge set of each circuit is closed under taking partners, and it
is $L$-connected (consecutive edges are $L$-adjacent). A nonempty
$L$-connected set closed under $L$-adjacency is the vertex set of one cycle
of $L$, so the partition coincides with the one just constructed.

\emph{Parity and length.} The dual lattice is bipartite: colour
$(a+\tfrac12,b+\tfrac12)$ by the parity of $a+b$; every step of a walk in
$\dualZ$ changes the colour, so every closed walk has even length. Hence
$r$ is even, and since $r\ge3$ we get $r\ge4$.
\end{proof}

\subsection{Contours as simple closed curves}

Let $\gamma^*$ be one of the circuits produced by
Lemma~\ref{lem:contour-decomp}. We realize it as a closed curve
$\widehat\gamma\subset\R^2$ as follows. Start from the union of the (closed)
dual edges of $\gamma^*$. At every passage of the circuit through a dual
vertex $x^*$ along two \emph{perpendicular} edges (a corner), delete from
those two edges the half-open segments of length $\tfrac14$ abutting $x^*$
(each containing $x^*$ but not the point at distance $\tfrac14$ from it)
and insert instead the straight chord joining the two surviving endpoints,
i.e.\ the two points at distance $\tfrac14$ from $x^*$ on the two edges
(``rounding the corner''). Straight passages are left untouched. Each chord
lies within distance $\tfrac14$ of its vertex $x^*$, inside the closed
quadrant at $x^*$ spanned by the directions of the two paired edges, and
avoids the point $x^*$ itself.

We claim that $\widehat\gamma$ is a \emph{simple} closed curve and that the
curves realizing distinct contours of the same configuration are pairwise
disjoint. Distinct dual edges are disjoint except possibly for a common
endpoint, and a chord at $x^*$ is disjoint from every dual edge not
containing $x^*$ (such an edge is at distance $\ge\tfrac12$ from $x^*$), so
it suffices to examine small neighbourhoods of the dual vertices. At a
vertex of degree $2$ in $\dedges(\omega)$ there is a single passage (of a
single contour) and nothing to check. At a vertex $x^*$ of degree $4$, all
four incident edges are trimmed and the two passages are the north--east
and the south--west corners; the two chords lie in the closed north-east,
respectively south-west, quadrant at $x^*$ and avoid $x^*$, so they are
disjoint from each other, and each meets the trimmed edges only at its two
endpoints, where the curve is glued. Hence no two passages intersect,
whether they belong to the same contour or to different ones, and the
claim follows.

By the Jordan curve theorem for piecewise-linear curves --- an elementary
classical result that we use without proof ---
$\R^2\setminus\widehat\gamma$ has exactly two connected components, one
bounded and one unbounded, each with boundary $\widehat\gamma$. This
theorem, together with the crossing-parity fact quoted in the proof of
Lemma~\ref{lem:edge-boundary} below, is the only topological input of this
section; both are standard, and they are the only statements below that are
sketched rather than proved in detail.

Note also that no site of $\Z^2$ lies on any realized curve: a site is at
distance $\tfrac12$ from every dual edge and at distance $\tfrac{\sqrt2}2$
from every dual vertex, hence at distance at least
$\tfrac{\sqrt2}2-\tfrac14>0$ from every chord. Thus every site of $\Z^2$
lies either in the bounded or in the unbounded component of
$\R^2\setminus\widehat\gamma$.

Finally, observe that the circuit structure, hence the curve
$\widehat\gamma$, is determined by the \emph{edge set} $\gamma^*$ alone,
independently of the ambient configuration: at every dual vertex, $\gamma^*$
contains either both edges of a single passage or all four incident edges
(if $e^*\in\gamma^*$ then so is its partner at each endpoint), and applying
the deformation rule to $\gamma^*$ itself reproduces the pairing inherited
from $\dedges(\omega)$. This makes the following definition unambiguous.

\begin{definition}\label{def:contours}
Let $\omega\in\Omega^{+}_{B(n)}$. The edge sets
$\gamma_1^*,\dots,\gamma_k^*$ of the unique decomposition of
$\dedges(\omega)$ given by Lemma~\ref{lem:contour-decomp} are called the
\emph{contours} of $\omega$, and we set
$\Gam(\omega)=\{\gamma_1^*,\dots,\gamma_k^*\}$ (with
$\Gam(\omega)=\varnothing$ when $\dedges(\omega)=\varnothing$). For a
contour $\gamma^*$, realized as the simple closed curve $\widehat\gamma$:
\begin{itemize}
  \item $|\gamma^*|$ denotes the number of dual edges of $\gamma^*$ (its
        \emph{length});
  \item $\Intr(\gamma^*)$ is the set of sites of $\Z^2$ lying in the
        bounded component of $\R^2\setminus\widehat\gamma$, and
        $\Extr(\gamma^*)$ the set of sites lying in the unbounded
        component, so that
        $\Z^2=\Intr(\gamma^*)\sqcup\Extr(\gamma^*)$; we say that $\gamma^*$
        \emph{surrounds} $i$ if $i\in\Intr(\gamma^*)$.
\end{itemize}
More generally, a set $\gamma^*$ of dual edges is called a \emph{contour in
$B(n)$} if $\gamma^*\in\Gam(\omega')$ for at least one
$\omega'\in\Omega^{+}_{B(n)}$, and for such $\gamma^*$ we write
\[
  \{\Gam\ni\gamma^*\}
  \;=\;\big\{\omega\in\Omega^{+}_{B(n)} : \gamma^*\in\Gam(\omega)\big\}
\]
for the event that $\gamma^*$ belongs to the contour family of the
configuration.
\end{definition}

\subsection{Three geometric lemmas}

The first lemma is the discrete Jordan-type fact on which everything else
rests: the dual edges of a contour are exactly the dual bonds separating its
interior from its exterior.

\begin{lemma}\label{lem:edge-boundary}
Let $\gamma^*$ be a contour in $B(n)$. Then:
\begin{enumerate}
  \item[(a)] for every nearest-neighbour edge $e=\{i,j\}\in\Eset_{\Z^2}$,
  \[
    e^*\in\gamma^*
    \quad\Longleftrightarrow\quad
    \text{exactly one of $i,j$ belongs to $\Intr(\gamma^*)$;}
  \]
  \item[(b)] $\Intr(\gamma^*)\subseteq B(n)$.
\end{enumerate}
\end{lemma}

\begin{proof}
(a) We use the following standard consequence of the Jordan curve theorem
(\emph{crossing parity}), which --- as announced --- we only sketch. Let
$p,q\in\R^2\setminus\widehat\gamma$ and suppose the segment $[p,q]$ meets
$\widehat\gamma$ in finitely many points, at each of which $\widehat\gamma$
is locally a straight segment crossed transversally by $[p,q]$. Then $p$ and
$q$ lie in the same component of $\R^2\setminus\widehat\gamma$ if and only
if the number of intersection points is even. \emph{Sketch:} on each maximal
subinterval of $[p,q]\setminus\widehat\gamma$ the ambient component is
constant; at a transversal crossing the segment passes from one local side
of the curve to the other; and for a simple closed curve the two local sides
of the curve near any of its points lie in the two distinct global
components. This last assertion is the standard refinement of the Jordan
curve theorem for polygonal curves, and is the step we do not prove.

Now fix $e=\{i,j\}$; by symmetry (the construction is invariant under the
lattice symmetries exchanging the axes) we may assume $e$ horizontal,
$i=(a,b)$, $j=(a+1,b)$. We count the points of $[i,j]\cap\widehat\gamma$:
\begin{itemize}
  \item every point of $[i,j]$ has second coordinate $b\in\Z$, while dual
        vertices have both coordinates in $\Z+\tfrac12$; hence $[i,j]$ is at
        distance $\ge\tfrac12$ from every dual vertex, and in particular
        avoids all chords (which lie within distance $\tfrac14$ of their
        vertices);
  \item a horizontal dual edge lies on a line $\{y=d+\tfrac12\}$, $d\in\Z$,
        hence is disjoint from $[i,j]$;
  \item a vertical dual edge
        $f^*=\{(c+\tfrac12,d+\tfrac12),(c+\tfrac12,d+\tfrac32)\}$ meets the
        line $\{y=b\}$ if and only if $b=d+1$, and then in the single point
        $(c+\tfrac12,b)$, which belongs to $[i,j]$ if and only if $c=a$.
        Thus the only dual edge meeting $[i,j]$ is $e^*$ itself, at its
        midpoint (which is also the midpoint of $e$);
  \item the midpoint of $e^*$ is at distance $\tfrac12$ from both endpoints
        of $e^*$, so it survives the trimming (which removes only points
        within distance $\tfrac14$ of dual vertices); near it,
        $\widehat\gamma$ coincides with the vertical segment $e^*$, which
        $[i,j]$ crosses transversally.
\end{itemize}
Hence $[i,j]$ meets $\widehat\gamma$ in exactly one point, a transversal
crossing, if $e^*\in\gamma^*$, and in no point otherwise. Since neither
$i$ nor $j$ lies on $\widehat\gamma$, the crossing-parity fact shows that
$i$ and $j$ lie in different components of $\R^2\setminus\widehat\gamma$ if
and only if $e^*\in\gamma^*$; and lying in different components means
precisely that exactly one of $i,j$ lies in the bounded component, i.e.\ in
$\Intr(\gamma^*)$.

(b) Let $K=[-n-\tfrac12,\,n+\tfrac12]^2$ and let
$\omega'\in\Omega^{+}_{B(n)}$ be a configuration with
$\gamma^*\in\Gam(\omega')$. Every $e^*\in\gamma^*\subseteq\dedges(\omega')$
is dual to a disagreement edge $e=\{i,j\}$, which has an endpoint of spin
$-1$, necessarily in $B(n)$; say $i\in B(n)$. In the direction of $e$, all
points of $e^*$ have the same coordinate, equal to the average of the
corresponding coordinates of $i$ and $j$; since the coordinate of $i$ lies
in $[-n,n]$ and that of $j$ differs from it by $1$, this average lies in
$[-n-\tfrac12,\,n+\tfrac12]$. In the perpendicular direction, the
coordinates of the points of $e^*$ range over $[c-\tfrac12,\,c+\tfrac12]$,
where $c\in[-n,n]$ is the common perpendicular coordinate of $i$ and $j$.
Hence $e^*\subseteq K$. Every chord of $\widehat\gamma$ is a straight
segment with both endpoints on dual edges of $\gamma^*$, hence in $K$, and
$K$ is convex; therefore $\widehat\gamma\subseteq K$. The set
$\R^2\setminus K$ is connected, unbounded and disjoint from
$\widehat\gamma$, so it is contained in a single component of
$\R^2\setminus\widehat\gamma$, necessarily the unbounded one. Consequently
the bounded component is contained in $K$, and
\[
  \Intr(\gamma^*)\;\subseteq\;K\cap\Z^2\;=\;\{-n,\dots,n\}^2\;=\;B(n).
  \qedhere
\]
\end{proof}

\begin{lemma}\label{lem:surround}
Let $\omega\in\Omega^{+}_{B(n)}$. If $\sg_0(\omega)=-1$, then there exists
$\gamma^*\in\Gam(\omega)$ with $0\in\Intr(\gamma^*)$.
\end{lemma}

\begin{proof}
For $i\in\Z^2$, let
\[
  N(i)\;=\;\#\big\{\gamma^*\in\Gam(\omega) : i\in\Intr(\gamma^*)\big\}.
\]
We prove the stronger identity
\[
  \sg_i(\omega)\;=\;(-1)^{N(i)}
  \qquad\text{for every } i\in\Z^2,
\]
which immediately yields the lemma: $\sg_0(\omega)=-1$ forces $N(0)$ to be
odd, in particular $N(0)\ge1$, so some contour of $\omega$ surrounds $0$.

Set $g(i)=\sg_i(\omega)\,(-1)^{N(i)}\in\{-1,+1\}$ and let
$\{i,j\}\in\Eset_{\Z^2}$ be arbitrary. For each contour
$\gamma^*\in\Gam(\omega)$, Lemma~\ref{lem:edge-boundary}(a) states that the
indicators $\mathbf 1_{\Intr(\gamma^*)}(i)$ and
$\mathbf 1_{\Intr(\gamma^*)}(j)$ differ if and only if
$\{i,j\}^*\in\gamma^*$, i.e.
\[
  \mathbf 1_{\Intr(\gamma^*)}(i)-\mathbf 1_{\Intr(\gamma^*)}(j)
  \;\equiv\;
  \mathbf 1\big\{\{i,j\}^*\in\gamma^*\big\}
  \pmod 2 .
\]
Summing over $\gamma^*\in\Gam(\omega)$, and using that the contours
partition $\dedges(\omega)$ (Lemma~\ref{lem:contour-decomp}), so that
$\{i,j\}^*$ belongs to exactly one contour if
$\{i,j\}^*\in\dedges(\omega)$ and to none otherwise, we obtain
\[
  N(i)-N(j)\;\equiv\;\mathbf 1\big\{\{i,j\}^*\in\dedges(\omega)\big\}
  \pmod 2 .
\]
On the other hand, by the definition of $\dedges(\omega)$,
\[
  \sg_i(\omega)\,\sg_j(\omega)
  \;=\;(-1)^{\mathbf 1\{\{i,j\}^*\in\dedges(\omega)\}} .
\]
Combining the two displays (and $N(i)+N(j)\equiv N(i)-N(j)\bmod 2$),
\[
  g(i)\,g(j)
  \;=\;\sg_i\sg_j\,(-1)^{N(i)+N(j)}
  \;=\;(-1)^{\mathbf 1\{\{i,j\}^*\in\dedges(\omega)\}}\,
       (-1)^{\mathbf 1\{\{i,j\}^*\in\dedges(\omega)\}}
  \;=\;+1 ,
\]
so $g(i)=g(j)$ for every nearest-neighbour pair: since $\Z^2$ is connected
under nearest-neighbour adjacency, $g$ is constant. Finally, pick any
$i\notin B(n)$: then $\sg_i(\omega)=+1$, and $N(i)=0$ because
$\Intr(\gamma^*)\subseteq B(n)$ for every contour
(Lemma~\ref{lem:edge-boundary}(b)). Hence $g\equiv+1$, which is the claimed
identity.
\end{proof}

\begin{lemma}\label{lem:cross}
Let $\gamma^*$ be a contour in $B(n)$ with $0\in\Intr(\gamma^*)$ and
$|\gamma^*|=\ell$. Then $\gamma^*$ contains a dual edge of the form
\[
  \big\{(k+\tfrac12,-\tfrac12),\,(k+\tfrac12,\tfrac12)\big\}
  \;=\;\big\{(k,0),(k+1,0)\big\}^*
  \qquad\text{for some integer } 0\le k\le \tfrac{\ell}{2}-1 .
\]
\end{lemma}

\begin{proof}
By Lemma~\ref{lem:edge-boundary}(b), $\Intr(\gamma^*)\subseteq B(n)$ is
finite, while it contains $0$ by hypothesis. Hence
\[
  k=\max\{m\ge0 : (m,0)\in\Intr(\gamma^*)\}
  \qquad\text{and}\qquad
  k'=\max\{m\ge0 : (-m,0)\in\Intr(\gamma^*)\}
\]
are well defined, and $k,k'\ge0$. The edge $\{(k,0),(k+1,0)\}$ has exactly
one endpoint in $\Intr(\gamma^*)$, so by Lemma~\ref{lem:edge-boundary}(a)
its dual edge $\{(k+\tfrac12,-\tfrac12),(k+\tfrac12,\tfrac12)\}$ belongs to
$\gamma^*$; likewise
$\{(-k'-\tfrac12,-\tfrac12),(-k'-\tfrac12,\tfrac12)\}
=\{(-k'-1,0),(-k',0)\}^*\in\gamma^*$.

It remains to bound $k$. Regard $\gamma^*$ as a closed circuit
$x_0^*,x_1^*,\dots,x_\ell^*=x_0^*$ in $\dualZ$
(Lemma~\ref{lem:contour-decomp}). The dual vertices
$u^*=(k+\tfrac12,\tfrac12)$ and $v^*=(-k'-\tfrac12,\tfrac12)$ are endpoints
of edges of $\gamma^*$, hence are visited by the circuit, and $u^*\neq v^*$
since their first coordinates satisfy $k+\tfrac12\ge\tfrac12$ and
$-k'-\tfrac12\le-\tfrac12$. The circuit therefore splits at $u^*$ and $v^*$
into two walks from $u^*$ to $v^*$, of lengths $\ell_1,\ell_2$ with
$\ell_1+\ell_2=\ell$. Every step of a walk in $\dualZ$ changes the first
coordinate by at most $1$, so each of the two walks has at least
\[
  \big|(k+\tfrac12)-(-k'-\tfrac12)\big| \;=\; k+k'+1 \;\ge\; k+1
\]
steps. Hence $\ell=\ell_1+\ell_2\ge 2(k+1)$, that is,
$k\le\tfrac{\ell}{2}-1$.
\end{proof}

\begin{remark}\label{rem:bijection}
The identity $\sg_i(\omega)=(-1)^{N(i)}$ established in the proof of
Lemma~\ref{lem:surround} shows that a configuration
$\omega\in\Omega^{+}_{B(n)}$ is reconstructible from its contour family, so
the map $\omega\mapsto\Gam(\omega)$ is \emph{injective}. With the
appropriate notion of admissibility (finite families of contours in $B(n)$
with pairwise disjoint edge sets whose union again satisfies the
even-degree property, with the deformation rule reproducing each member),
the map is in fact a \emph{bijection} onto the set of admissible families of
pairwise compatible contours: reading the identity backwards produces, from
any such family, the unique configuration admitting it as its contour
family. We will not need this surjectivity statement: the Peierls argument
in Sections~\ref{sec:peierls}--\ref{sec:conclusion} only uses the events
$\{\Gam\ni\gamma^*\}$ together with Lemmas~\ref{lem:surround}
and~\ref{lem:cross}, so we omit the (easy) proof.
\end{remark}

\section{Peierls' estimate}
\label{sec:peierls}

\begin{lemma}[Peierls' estimate; Lemma 3.37, eq.~(3.34) of the transcription]
\label{lem:peierls}
For all $\beta>0$, all $n$, and any contour $\gamma^*$,
\[
  \mu^{+}_{B(n);\beta,0}(\Gam\ni\gamma^*)\le e^{-2\beta|\gamma^*|}.
\]
\end{lemma}

\noindent
The proof follows the classical energy argument: erasing the contour
$\gamma^*$ from a configuration lowers the energy by exactly
$2\beta|\gamma^*|$, and the erasure map is injective, so the constrained
partition sum is exponentially small compared with the full one. We mark the
steps explicitly; everything rests on
Lemmas~\ref{lem:edge-boundary} and~\ref{lem:contour-decomp}.

\begin{proof}
Fix $n$, $\beta>0$ and a contour $\gamma^*$. Throughout, write
$\mu^{+}=\mu^{+}_{B(n);\beta,0}$, $\Ham=\Ham^{+}_{B(n);\beta,0}$,
$Z=Z^{+}_{B(n);\beta,0}$, and let
\[
  \Eset_{B(n)}
  =\bigl\{\{i,j\}\in\Eset_{\Z^2} : \{i,j\}\cap B(n)\neq\varnothing\bigr\}
\]
denote the (finite) set of edges appearing in the Hamiltonian of
Section~\ref{sec:setting}.

\medskip\noindent\emph{Step 1: the erasure map.}
Define $T_{\gamma^*}:\Omega^{+}_{B(n)}\to\{-1,1\}^{\Z^2}$ by
\[
  (T_{\gamma^*}\omega)_i=
  \begin{cases}
    -\omega_i & \text{if } i\in\Intr(\gamma^*),\\
    \ \ \omega_i & \text{otherwise.}
  \end{cases}
\]
By Lemma~\ref{lem:edge-boundary}(b), $\Intr(\gamma^*)\subseteq B(n)$, so
$T_{\gamma^*}$ does not modify any spin outside $B(n)$; hence
$T_{\gamma^*}\omega\in\Omega^{+}_{B(n)}$, i.e.\ $T_{\gamma^*}$ maps
$\Omega^{+}_{B(n)}$ into itself. Moreover $T_{\gamma^*}\circ T_{\gamma^*}$ is
the identity, so $T_{\gamma^*}$ is an involution of $\Omega^{+}_{B(n)}$; in
particular it is injective, and a fortiori injective when restricted to the
event $\{\Gam\ni\gamma^*\}=\{\omega:\gamma^*\in\Gam(\omega)\}$.

\medskip\noindent\emph{Step 2: the energy identity.}
We claim that for every $\omega\in\Omega^{+}_{B(n)}$ with
$\gamma^*\in\Gam(\omega)$,
\begin{equation}
  \label{eq:energy-drop}
  \Ham(T_{\gamma^*}\omega)=\Ham(\omega)-2\beta|\gamma^*| .
\end{equation}
Both sides are finite sums over the same edge set $\Eset_{B(n)}$, so it
suffices to compare the two Hamiltonians edge by edge. Fix
$\{i,j\}\in\Eset_{\Z^2}$ and consider the product
$\sg_i\sg_j$. Since $T_{\gamma^*}$ flips precisely the spins in
$\Intr(\gamma^*)$,
\[
  \sg_i(T_{\gamma^*}\omega)\,\sg_j(T_{\gamma^*}\omega)
  =-\,\sg_i(\omega)\,\sg_j(\omega)
  \quad\Longleftrightarrow\quad
  \text{exactly one of } i,j \text{ lies in } \Intr(\gamma^*),
\]
and the product is unchanged when both or neither endpoint lies in
$\Intr(\gamma^*)$. By Lemma~\ref{lem:edge-boundary}(a), the edges with
exactly one endpoint in $\Intr(\gamma^*)$ are exactly the $|\gamma^*|$ edges
of $\Eset_{\Z^2}$ whose dual edge belongs to $\gamma^*$.

Now use the hypothesis $\gamma^*\in\Gam(\omega)$. By
Lemma~\ref{lem:contour-decomp}, the edges of $\gamma^*$ all belong to
$\dedges(\omega)$; that is, each edge $\{i,j\}$ dual to an edge of $\gamma^*$
is a disagreement edge of $\omega$:
$\sg_i(\omega)\sg_j(\omega)=-1$. Two consequences:
\begin{enumerate}
\item Every such $\{i,j\}$ lies in $\Eset_{B(n)}$, and so genuinely appears
  in the Hamiltonian. Indeed, if both endpoints lay outside $B(n)$, then
  $\omega\in\Omega^{+}_{B(n)}$ would force
  $\sg_i(\omega)=\sg_j(\omega)=+1$ and the edge could not be a disagreement
  edge. (Alternatively: by Lemma~\ref{lem:edge-boundary} one endpoint lies in
  $\Intr(\gamma^*)\subseteq B(n)$.)
\item With the sign convention
  $\Ham(\omega)=-\beta\sum_{\{i,j\}\in\Eset_{B(n)}}\sg_i(\omega)\sg_j(\omega)$,
  each such edge contributes $-\beta\cdot(-1)=+\beta$ to $\Ham(\omega)$,
  while after the flip it satisfies
  $\sg_i(T_{\gamma^*}\omega)\sg_j(T_{\gamma^*}\omega)=+1$ and contributes
  $-\beta$ to $\Ham(T_{\gamma^*}\omega)$: a decrease of $2\beta$ per edge of
  $\gamma^*$.
\end{enumerate}
All remaining edges of $\Eset_{B(n)}$ have zero or two endpoints in
$\Intr(\gamma^*)$ and contribute equally to $\Ham(\omega)$ and to
$\Ham(T_{\gamma^*}\omega)$. Summing over $\Eset_{B(n)}$ yields
\eqref{eq:energy-drop}: erasing the contour lowers the energy by exactly
$2\beta|\gamma^*|$.

\medskip\noindent\emph{Step 3: conclusion.}
Exponentiating \eqref{eq:energy-drop}, for every $\omega$ with
$\gamma^*\in\Gam(\omega)$ we have
$e^{-\Ham(\omega)}=e^{-2\beta|\gamma^*|}\,e^{-\Ham(T_{\gamma^*}\omega)}$.
Therefore
\[
  \mu^{+}(\Gam\ni\gamma^*)
  =\frac{1}{Z}\sum_{\substack{\omega\in\Omega^{+}_{B(n)}:\\ \Gam(\omega)\ni\gamma^*}}
     e^{-\Ham(\omega)}
  =e^{-2\beta|\gamma^*|}\,
   \frac{1}{Z}\sum_{\substack{\omega\in\Omega^{+}_{B(n)}:\\ \Gam(\omega)\ni\gamma^*}}
     e^{-\Ham(T_{\gamma^*}\omega)} .
\]
Since $T_{\gamma^*}$ is injective on $\{\Gam\ni\gamma^*\}$ (Step 1), the
configurations $T_{\gamma^*}\omega$ appearing in the last sum are pairwise
distinct elements of $\Omega^{+}_{B(n)}$, so the sum of their Boltzmann
weights is bounded by the full partition function:
\[
  \sum_{\substack{\omega\in\Omega^{+}_{B(n)}:\\ \Gam(\omega)\ni\gamma^*}}
     e^{-\Ham(T_{\gamma^*}\omega)}
  \;\le\;
  \sum_{\omega'\in\Omega^{+}_{B(n)}} e^{-\Ham(\omega')}
  \;=\; Z .
\]
Combining the two displays gives
$\mu^{+}(\Gam\ni\gamma^*)\le e^{-2\beta|\gamma^*|}$, as claimed. Note that
the bound is uniform in $n$ and holds for every $\beta>0$.
\end{proof}

\begin{remark}
As the transcription observes after Lemma~3.37, the estimate expresses the
\emph{stability of the ground state} $\eta^+\equiv+1$ at large $\beta$: any
fixed deviation from $\eta^+$, encoded by a contour $\gamma^*$, has
probability exponentially small in the energy cost $2\beta|\gamma^*|$ of its
boundary. (Beyond the book, our phrasing:) the argument trades energy against
entropy --- the exponential decay in $|\gamma^*|$ obtained here will beat, for
$\beta$ large, the exponential growth in the \emph{number} of contours of a
given length established in Section~\ref{sec:counting}.
\end{remark}

\section{Counting contours}
\label{sec:counting}

The second ingredient of the Peierls argument is an \emph{entropy} bound: an
upper bound on the number of contours of a given length, first through a fixed
dual edge, then surrounding the origin. Together with the \emph{energy} bound
of Lemma~\ref{lem:peierls} it will make the sum over contours convergent and
small at large $\beta$. Throughout, ``contour'' means a contour in $B(n)$ for
some (arbitrary) $n$, in the sense of Definition~\ref{def:contours}; the
bounds below use only the circuit property of
Lemma~\ref{lem:contour-decomp} and are uniform in $n$.

\begin{lemma}
\label{lem:walkcount}
Let $e^*$ be a fixed edge of the dual lattice $\dualZ$ and let $\ell\ge1$.
The number of contours $\gamma^*$ with $|\gamma^*|=\ell$ and
$e^*\in\gamma^*$ is at most $3^{\ell-1}$.
\end{lemma}

\begin{proof}
By Lemma~\ref{lem:contour-decomp}, every contour is a closed
edge-self-avoiding circuit in $\dualZ$. Fix one of the two endpoints of $e^*$
as starting vertex, i.e.\ fix once and for all a direction in which $e^*$ is
crossed. Every contour $\gamma^*$ of length $\ell$ containing $e^*$ can then
be traversed as a closed walk
\[
  e^*_1=e^*,\ e^*_2,\ \dots,\ e^*_\ell
\]
in $\dualZ$ which starts by crossing $e^*$ in the fixed direction and uses
each of the $\ell$ edges of $\gamma^*$ exactly once: one starts the
traversal of the circuit at an occurrence of $e^*$, reversing the circuit's
orientation if necessary so that $e^*$ is crossed in the fixed direction
(reversal is legitimate, since the pairing condition of
Lemma~\ref{lem:contour-decomp} involves unordered pairs of edges). We count
such walks. The first step
is the prescribed crossing of $e^*$, so it offers no choice. After
$k$ steps, $1\le k\le\ell-1$, the walk sits at a dual vertex, which has degree
$4$ in $\dualZ$; the next edge $e^*_{k+1}$ is one of the $4$ edges incident
to that vertex, but it cannot be the edge $e^*_k$ just traversed, since
immediately reversing would use the same edge twice, contradicting
edge-self-avoidance. Hence each of the $\ell-1$ steps after the first admits
at most $3$ choices, and there are at most $3^{\ell-1}$ such walks. Since
every contour of length $\ell$ containing $e^*$ arises (as the set of edges)
from at least one of these walks, the number of such contours is at most
$3^{\ell-1}$.
\end{proof}

\begin{remark}
\emph{(Our remark, not part of the source text.)} The transcription records
as Lemma~3.38 the Eulerian-traversal counting idea: a connected graph with
$N$ edges admits a closed walk using each edge at most twice, which bounds
the number of admissible edge sets of size $N$ through a fixed edge by an
exponential in $N$. Because our contours are already closed
edge-self-avoiding \emph{circuits} (Lemma~\ref{lem:contour-decomp}), the
simpler one-pass walk count of Lemma~\ref{lem:walkcount} suffices here; it is
a variant of the same idea, traversing each edge once instead of at most
twice.
\end{remark}

\begin{lemma}
\label{lem:counting}
For every even $\ell\ge4$, the number of contours $\gamma^*$ with
$|\gamma^*|=\ell$ and $0\in\Intr(\gamma^*)$ is at most
$\dfrac{\ell}{2}\,3^{\ell-1}$.
\end{lemma}

\begin{proof}
By Lemma~\ref{lem:cross}, every contour of length $\ell$ surrounding the
origin contains at least one of the $\ell/2$ vertical dual edges
\[
  e^*_k=\big\{(k+\tfrac12,-\tfrac12),\,(k+\tfrac12,\tfrac12)\big\},
  \qquad k\in\{0,1,\dots,\tfrac{\ell}{2}-1\}.
\]
For each fixed $k$, Lemma~\ref{lem:walkcount} bounds the number of contours
of length $\ell$ containing $e^*_k$ by $3^{\ell-1}$. Summing over the
$\ell/2$ possible values of $k$ yields the claim.
\end{proof}

\section{Proof of the main theorem}
\label{sec:conclusion}

We now assemble the geometric input (Lemmas~\ref{lem:contour-decomp},
\ref{lem:surround} and \ref{lem:cross}), the energy estimate
(Lemma~\ref{lem:peierls}) and the entropy estimate
(Lemma~\ref{lem:counting}) into a proof of the main theorem.

\begin{proof}[Proof of Theorem~\ref{thm:main}]
We show that the choice $\beta_0=1$ works. Fix $n\ge1$ and $\beta\ge1$, and
abbreviate $\mu=\mu^{+}_{B(n);\beta,0}$.

\medskip\noindent
\emph{Step 1: finite-volume bound on $\mu(\sg_0=-1)$, uniformly in $n$.}
By Lemma~\ref{lem:surround}, if $\sg_0(\omega)=-1$ then some contour of
$\Gam(\omega)$ surrounds the origin, i.e.
\[
  \{\sg_0=-1\}\ \subset
  \bigcup_{\gamma^*:\,0\in\Intr(\gamma^*)}\{\Gam\ni\gamma^*\},
\]
where the union runs over the set of \emph{all} contours surrounding the
origin --- not merely those appearing in a given configuration. This set is
countable: by Lemma~\ref{lem:contour-decomp} every contour has even length
$\ell\ge4$, and by Lemma~\ref{lem:counting} only finitely many contours of
each length surround the origin. Hence the union bound gives a well-defined
sum, which we organize by length and estimate by Peierls' estimate
(Lemma~\ref{lem:peierls}) and Lemma~\ref{lem:counting}:
\[
  \mu(\sg_0=-1)
  \ \le \sum_{\gamma^*:\,0\in\Intr(\gamma^*)}\mu(\Gam\ni\gamma^*)
  \ \le \sum_{\substack{\ell\ge4\\ \ell\ \mathrm{even}}}
        \frac{\ell}{2}\,3^{\ell-1}\,e^{-2\beta\ell}.
\]
Since $\frac{\ell}{2}\,3^{\ell-1}=\frac{\ell}{6}\,3^{\ell}$, and since
extending the sum from even $\ell\ge4$ to all integers $\ell\ge4$ only adds
nonnegative terms, we obtain, with
\[
  x=x(\beta)=3e^{-2\beta},
\]
the bound
\[
  \mu^{+}_{B(n);\beta,0}(\sg_0=-1)\ \le\ \frac16\sum_{\ell\ge4}\ell\,x^{\ell}.
\]
The right-hand side does not depend on $n$: the bound is uniform in the
volume.

\medskip\noindent
\emph{Step 2: evaluation of the series.} For $0\le x<1$ and any integer
$N\ge1$, all the series below converge absolutely, and
\[
  (1-x)\sum_{\ell\ge N}\ell\,x^{\ell}
  =\sum_{\ell\ge N}\ell\,x^{\ell}-\sum_{\ell\ge N+1}(\ell-1)\,x^{\ell}
  =N x^{N}+\sum_{\ell\ge N+1}x^{\ell}
  =N x^{N}+\frac{x^{N+1}}{1-x},
\]
whence
\[
  \sum_{\ell\ge N}\ell\,x^{\ell}
  =\frac{x^{N}\big(N-(N-1)x\big)}{(1-x)^{2}},
  \qquad\text{in particular}\qquad
  \sum_{\ell\ge4}\ell\,x^{\ell}=\frac{x^{4}(4-3x)}{(1-x)^{2}}.
\]
(As a sanity check at $N=2$, $x=\tfrac12$: the formula gives
$\tfrac14\cdot\tfrac32\,/\,\tfrac14=\tfrac32=\sum_{\ell\ge2}\ell\,2^{-\ell}$.)
Combining with Step 1,
\[
  \mu^{+}_{B(n);\beta,0}(\sg_0=-1)\ \le\ f\big(x(\beta)\big),
  \qquad
  f(x)=\frac{x^{4}(4-3x)}{6(1-x)^{2}}.
\]

\medskip\noindent
\emph{Step 3: numerical bound, valid for all $\beta\ge1$.} Since
$e^{2}>7.38$, we have
\[
  x(1)=3e^{-2}<\frac{3}{7.38}<0.41,
\]
and $\beta\mapsto x(\beta)=3e^{-2\beta}$ is decreasing, so
$0<x(\beta)<0.41$ for every $\beta\ge1$. Moreover $f$ is increasing on
$[0,0.41]$: the numerator $x^{4}(4-3x)=4x^{4}-3x^{5}$ has derivative
$16x^{3}-15x^{4}=x^{3}(16-15x)>0$ for $0<x\le1$, and $(1-x)^{-2}$ is positive
and increasing on $[0,1)$. Hence $f(x(\beta))\le f(0.41)$, which we bound by
elementary inequalities:
\[
  (0.41)^{4}=0.02825761<0.0283,\qquad
  4-3\cdot(0.41)<4,\qquad
  (1-0.41)^{2}=(0.59)^{2}=0.3481,
\]
so that
\[
  f(0.41)\ <\ \frac{0.0283\cdot 4}{6\cdot 0.3481}
          \ =\ \frac{0.1132}{2.0886}\ <\ 0.0542\ <\ \frac14
\]
(indeed $0.0542\cdot2.0886=0.11320212>0.1132$). Therefore
\[
  \mu^{+}_{B(n);\beta,0}(\sg_0=-1)\ \le\ 0.0542\ <\ \tfrac14
  \qquad\text{for all } n\ge1 \text{ and all } \beta\ge1 .
\]

\medskip\noindent
\emph{Step 4: finite-volume magnetization.} On $\Omega^{+}_{B(n)}$ the spin
$\sg_0$ takes values in $\{-1,1\}$, and
$\sg_0=1-2\cdot\mathbf{1}_{\{\sg_0=-1\}}$: both sides equal $1$ when
$\sg_0=1$ and $-1$ when $\sg_0=-1$. Taking expectations and using Step 3,
\[
  \avg{\sg_0}^{+}_{B(n);\beta,0}
  =1-2\,\mu^{+}_{B(n);\beta,0}(\sg_0=-1)
  \ \ge\ 1-2\cdot\tfrac14=\tfrac12
  \qquad\text{for all } n\ge1,\ \beta\ge1 .
\]

\medskip\noindent
\emph{Step 5: thermodynamic limit.} By Theorem~3.17 of the transcription, the
limit
$\meanp{\sg_0}_{\beta,0}=\lim_{n\to\infty}\avg{\sg_0}^{+}_{B(n);\beta,0}$
exists (and is independent of the sequence of volumes). Since each term of
the sequence is at least $\tfrac12$, so is the limit: in the notation of
Definition~3.32 of the transcription,
\[
  m^*(\beta)=\meanp{\sg_0}_{\beta,0}\ \ge\ \tfrac12\ >\ 0
  \qquad\text{for every }\beta\ge1 .
\]

\medskip\noindent
\emph{Step 6: non-uniqueness and $\beta_c(2)<\infty$.} By Proposition~3.29 of
the transcription, $\meanm{\sg_0}_{\beta,0}=-\meanp{\sg_0}_{\beta,0}$, so for
every $\beta\ge1$
\[
  \meanm{\sg_0}_{\beta,0}\le-\tfrac12<\tfrac12\le\meanp{\sg_0}_{\beta,0}.
\]
Thus $\meanp{\cdot}_{\beta,0}$ and $\meanm{\cdot}_{\beta,0}$ are two
\emph{distinct} Gibbs states: this is non-uniqueness in the sense of
Definition~3.24 of the transcription (equivalently, via the one-point
criterion of Theorem~3.28(3), since
$\meanp{\sg_0}_{\beta,0}\neq\meanm{\sg_0}_{\beta,0}$). Finally, since
$m^*(\beta)>0$ for every $\beta\ge1$, Definition~3.32 gives
\[
  \beta_c(2)=\sup\{\beta\ge0:m^*(\beta)=0\}\ \le\ 1\ =\ \beta_0\ <\ \infty .
\]
This proves all assertions of Theorem~\ref{thm:main} with $\beta_0=1$.
\end{proof}

\begin{remark}
\label{rem:d3}
\emph{(Our remark; sketch only, beyond what is proved in this document.)}
The same conclusion extends to all dimensions $d\ge3$, with the bound
$\beta_c(d)\le\beta_c(2)<\infty$. Indeed, by a standard consequence of the
GKS inequalities (Theorem~3.20 of the transcription), at $h=0$ the expectations
$\meanp{\sg_A}$ are nondecreasing functions of the coupling constants
$J\ge0$. Embedding $\Z^d$ into $\Z^{d+1}$ as a coordinate hyperplane and
setting to zero the couplings on all edges leaving that hyperplane reduces
the $(d+1)$-dimensional model to the $d$-dimensional one; restoring those
couplings to their ferromagnetic value can, by GKS monotonicity, only
increase $\meanp{\sg_0}$. Hence $m^*_{d+1}(\beta)\ge m^*_{d}(\beta)$ for all
$\beta\ge0$, and $\beta_c(d+1)\le\beta_c(d)$. We do not carry out the details
here; the statement is consistent with the transcription's claim
(\S3.7.2 there) that $\beta_c(d)<\infty$ for all $d\ge2$.
\end{remark}

\begin{remark}
\label{rem:explicit}
\emph{(Our remark; the cited improvements are not proved in this document.)}
The value $\beta_0=1$ obtained above is far from optimal. A more careful
version of the contour counting --- recorded as Exercise~3.21 in the book,
according to the transcription (\S3.7.2 there) --- improves the conclusion to
the explicit bound $\beta_c(2)<0.88$. For perspective, the exact value in
dimension two is known from the Onsager solution
(see \S3.10.12 of the transcription):
$\beta_c(2)=\tfrac12\log(1+\sqrt2)\approx0.4407$. Neither refinement is
proved here.
\end{remark}

\end{document}
